% AA-20 | Algebra abstrakcyjna | Iteracja X
\chapter{Rozszerzenia ciał i ciała rozkładu wielomianów}
\label{aa:ch:field-extensions}

\section{Rozszerzenie ciała jako przestrzeń wektorowa}
Gdy rozważamy rozwiązania równań wielomianowych, nawet nad
ciałem może zabraknąć pierwiastków. Przykładowo
$X^2-2$ nie ma pierwiastka w $\mathbb Q$, choć ma go
w większym ciele. Chcemy systematycznie konstruować
większe ciała i mierzyć, ile nowych elementów trzeba dodać.

\begin{definicja}[Rozszerzenie ciał i jego stopień]
\label{aa:def:field-extension-degree}
Niech $K,L$ będą ciałami, a $K\subseteq L$ podciałem.
Mówimy, że $L/K$ jest \emph{rozszerzeniem ciał}.
Mnożenie w $L$ ograniczone do $K\times L$ nadaje $L$
strukturę przestrzeni wektorowej nad $K$.
Jej wymiar oznaczamy $[L:K]$ i nazywamy
\emph{stopniem rozszerzenia}. Rozszerzenie jest
\emph{skończone}, jeśli $[L:K]<\infty$.
W przypadku zanurzenia $K\hookrightarrow L$ utożsamiamy
$K$ z jego obrazem, aby móc stosować ten sam zapis.
\end{definicja}

\begin{przyklad}[Stopień rozszerzenia $\mathbb C/\mathbb R$]
\label{aa:prz:complex-degree}
Z konstrukcji $\mathbb C\cong\mathbb R^2$ w
Rozdziale~\ref{aa:ch:fields-fractions} wynika, że
każdy $z\in\mathbb C$ ma jednoznacznie postać
$z=a+bi$ z $a,b\in\mathbb R$. Wektory $1,i$
tworzą więc bazę $\mathbb C$ nad $\mathbb R$, a zatem
$[\mathbb C:\mathbb R]=2$. Ten fakt nie wymaga
pełnej teorii funkcji holomorficznych ani konstruowania
liczb rzeczywistych na nowo.
\end{przyklad}

\begin{twierdzenie}[Twierdzenie o wieży rozszerzeń]
\label{aa:tw:field-tower}
Niech $K\subseteq E\subseteq L$ będą ciałami.
Jeśli $[E:K]=r<\infty$ i $[L:E]=s<\infty$,
to $[L:K]=rs$. Ponadto, jeśli $[L:K]<\infty$,
to oba stopnie po prawej stronie są skończone.
\end{twierdzenie}
\begin{proof}
Wybierzmy bazę $u_1,\dots,u_r$ przestrzeni $E$ nad $K$
oraz bazę $v_1,\dots,v_s$ przestrzeni $L$ nad $E$.
Pokażemy, że $rs$ iloczynów $u_iv_j$ tworzy bazę $L$ nad $K$.

\emph{Rozpinanie.} Dla $x\in L$ istnieją $b_j\in E$ takie,
że $x=\sum_j b_jv_j$. Dla każdego $j$ istnieją $a_{ij}\in K$
z $b_j=\sum_i a_{ij}u_i$. Wobec tego
$x=\sum_{j=1}^s\sum_{i=1}^r a_{ij}(u_iv_j)$,
czyli wskazane iloczyny rozpinają $L$.

\emph{Niezależność.} Jeśli
$\sum_{j=1}^s\sum_{i=1}^r a_{ij}u_iv_j=0$,
to traktując $\sum_i a_{ij}u_i$ jako współczynnik w $E$
i korzystając z niezależności $v_j$ nad $E$ dostajemy
$\sum_i a_{ij}u_i=0$ dla każdego $j$.
Niezależność $u_i$ nad $K$ daje wszystkie $a_{ij}=0$.
Stąd $[L:K]=rs$.

Jeśli natomiast $[L:K]=n<\infty$, to $E$ jest
$K$-podprzestrzenią $L$, a zatem $[E:K]\le n$.
Każda rodzina liniowo niezależna elementów $L$ nad $E$
jest też niezależna nad $K$: każda relacja o współczynnikach
w $K$ byłaby również relacją nad $E$. Stąd $[L:E]\le n$.
Możemy więc zastosować udowodnioną pierwszą część.
\end{proof}

\begin{wniosek}[Podzielność stopni]
\label{aa:wn:tower-divisibility}
Jeśli $L/K$ jest rozszerzeniem skończonym i
$K\subseteq E\subseteq L$, to $[E:K]$ dzieli $[L:K]$.
\end{wniosek}
\begin{proof}
Równość $[L:K]=[L:E][E:K]$ przedstawia $[L:K]$
jako iloczyn dwóch dodatnich liczb całkowitych.
\end{proof}

\section{Elementy algebraiczne i wielomiany minimalne}
\begin{definicja}[Algebraiczność i przestępność]
\label{aa:def:algebraic-element}
Niech $L/K$ będzie rozszerzeniem ciał, a $\alpha\in L$.
Element $\alpha$ jest \emph{algebraiczny nad $K$}, jeśli
istnieje niezerowy wielomian $f\in K[X]$ z $f(\alpha)=0$.
W przeciwnym przypadku jest \emph{przestępny nad $K$}.
Przez $K[\alpha]$ oznaczamy najmniejszy podpierścień $L$
zawierający $K$ i $\alpha$, a przez $K(\alpha)$
najmniejsze podciało $L$ zawierające te elementy.
\end{definicja}

\begin{lemat}[Wielomian minimalny]
\label{aa:lem:minimal-polynomial}
Niech $L/K$ będzie rozszerzeniem ciał i niech
$\alpha\in L$ będzie algebraiczny nad $K$. Wówczas istnieje
jedyny wielomian unormowany $m_{\alpha,K}\in K[X]$
o najmniejszym dodatnim stopniu spełniający
$m_{\alpha,K}(\alpha)=0$. Wielomian ten jest
nierozkładalny i dzieli każdy $f\in K[X]$
taki, że $f(\alpha)=0$.
\end{lemat}
\begin{proof}
Rozważmy ideał
$I_\alpha=\{f\in K[X]:f(\alpha)=0\}$.
Jest to jądro homomorfizmu podstawiania
$\operatorname{ev}_\alpha:K[X]\to L$.
Ponieważ $\alpha$ jest algebraiczny, ideał zawiera
niezerowy wielomian. Nie zawiera jednak niezerowej
stałej: dla $a\in K$ mamy $a(\alpha)=a$, a $K$ jest
podciałem $L$. Wybierzmy w nim wielomian $m$
o najmniejszym możliwym stopniu, a następnie podzielmy
przez jego niezerowy współczynnik wiodący, aby był unormowany.

Dla dowolnego $f\in I_\alpha$ dzielenie z resztą
w $K[X]$ daje $f=qm+r$, gdzie $r=0$ albo
$\deg r<\deg m$. Po podstawieniu $\alpha$
otrzymujemy $r(\alpha)=f(\alpha)-q(\alpha)m(\alpha)=0$.
Minimalność stopnia $m$ wymusza $r=0$, więc $m\mid f$.
Jeżeli $m=uv$ dla wielomianów dodatniego stopnia,
to $0=m(\alpha)=u(\alpha)v(\alpha)$.
Ponieważ $L$ jest ciałem, jeden z tych dwóch czynników
musi być zerowy. Byłby to wielomian mniejszego stopnia
niż $m$, sprzecznie z jego wyborem. Wielomian $m$ jest
zatem nierozkładalny.

Jeżeli $m'$ jest innym unormowanym wielomianem o żądanych
własnościach, to $m\mid m'$ i $m'\mid m$.
Ich stopnie są równe, a unormowanie wymusza $m=m'$.
\end{proof}

\begin{twierdzenie}[Proste rozszerzenie algebraiczne]
\label{aa:tw:simple-algebraic-extension}
Niech $L/K$ będzie rozszerzeniem ciał i
$\alpha\in L$ elementem algebraicznym nad $K$.
Jeśli $d=\deg m_{\alpha,K}$, to
\[
K[\alpha]=K(\alpha)\cong K[X]/(m_{\alpha,K}),
\qquad [K(\alpha):K]=d.
\]
Bazę $K(\alpha)$ nad $K$ stanowi rodzina
$1,\alpha,\ldots,\alpha^{d-1}$.
\end{twierdzenie}
\begin{proof}
Rozważmy homomorfizm
$\operatorname{ev}_\alpha:K[X]\to L$ z
$f\mapsto f(\alpha)$. Jego obrazem jest $K[\alpha]$,
a jądro jest równe $(m_{\alpha,K})$ na mocy
Lematu~\ref{aa:lem:minimal-polynomial}.
Pierwsze twierdzenie o izomorfizmie pierścieni daje
$K[X]/(m_{\alpha,K})\cong K[\alpha]$.

Dlaczego iloraz jest ciałem? Jeśli $[f]\ne0$, to
$m_{\alpha,K}\nmid f$. Nierozkładalność $m_{\alpha,K}$
i algorytm Euklidesa w $K[X]$ dają
$\gcd(f,m_{\alpha,K})=1$. Z tożsamości Bézouta
istnieją $u,v\in K[X]$ z
$uf+vm_{\alpha,K}=1$. W ilorazie mamy
$[u][f]=[1]$, więc każda niezerowa klasa ma odwrotność.
Stąd obraz $K[\alpha]$ jest już ciałem, a więc
$K[\alpha]=K(\alpha)$.

Dzielenie wielomianu przez $m_{\alpha,K}$ daje
jedyny reprezentant stopnia mniejszego niż $d$.
Dlatego każdy element ma postać
$\sum_{j=0}^{d-1}a_j\alpha^j$ z $a_j\in K$.
Jeśli taka suma jest zerowa, to odpowiedni wielomian
leży w $(m_{\alpha,K})$ i ma stopień mniejszy niż $d$;
musi więc być zerowy. Współczynniki są zerowe,
a wskazana rodzina jest bazą.
\end{proof}

\begin{przyklad}[Rozszerzenie przez pierwiastek z dwóch]
\label{aa:prz:sqrt-two-extension}
Wielomian $X^2-2$ jest nierozkładalny nad
$\mathbb Q$: gdyby miał wymierny pierwiastek $a/b$
w postaci nieskracalnej, to z $a^2=2b^2$ wynikałoby,
że $a$ jest parzyste, a potem także $b$, sprzecznie
z nieskracalnością. Dlatego
$[\mathbb Q(\sqrt2):\mathbb Q]=2$, a każdy element
ma jednoznacznie postać $a+b\sqrt2$ z $a,b\in\mathbb Q$.
Dla $a^2-2b^2\ne0$ odwrotność otrzymujemy przez
pomnożenie przez sprzężenie:
\[
(a+b\sqrt2)^{-1}=\frac{a-b\sqrt2}{a^2-2b^2}.
\]
Mianownik nie może być zerem dla $a+b\sqrt2\ne0$,
bo równość $a^2=2b^2$ przy $b\ne0$ dawałaby
wymierny pierwiastek z dwóch.
\end{przyklad}

\begin{przyklad}[Rozszerzenie o dwa pierwiastki kwadratowe]
\label{aa:prz:biquadratic-extension}
Niech $E=\mathbb Q(\sqrt2,\sqrt3)\subseteq\mathbb R$.
Wiemy, że $[\mathbb Q(\sqrt2):\mathbb Q]=2$.
Sprawdźmy, że $\sqrt3\notin\mathbb Q(\sqrt2)$.
Gdyby $\sqrt3=a+b\sqrt2$ z $a,b\in\mathbb Q$,
to po podniesieniu do kwadratu mielibyśmy
$3=a^2+2b^2+2ab\sqrt2$. Niewymierność $\sqrt2$
wymusza $ab=0$. Jeśli $b=0$, mielibyśmy
$3=a^2$ z $a\in\mathbb Q$, co jest niemożliwe:
wykładnik liczby pierwszej $3$ w rozkładzie kwadratu
liczby wymiernej jest parzysty, a w $3$ wynosi $1$.
Jeśli $a=0$, mielibyśmy $3=2b^2$, czyli
$b^2=3/2$, również niemożliwe, ponieważ
wykładnik liczby pierwszej $2$ w rozkładzie
kwadratu liczby wymiernej jest parzysty,
a w $3/2$ jest nieparzysty.
Wobec tego $X^2-3$ jest nierozkładalny nad
$\mathbb Q(\sqrt2)$ i
$[E:\mathbb Q(\sqrt2)]=2$.
Twierdzenie o wieży daje $[E:\mathbb Q]=4$,
a bazą jest $1,\sqrt2,\sqrt3,\sqrt6$.
\end{przyklad}

\begin{stwierdzenie}[Rozszerzenie skończone jest algebraiczne]
\label{aa:st:finite-implies-algebraic}
Niech $L/K$ będzie rozszerzeniem ciał stopnia
$[L:K]=d<\infty$. Wówczas każdy element $\alpha\in L$
jest algebraiczny nad $K$, a
$\deg m_{\alpha,K}\le d$.
\end{stwierdzenie}
\begin{proof}
Rodzina $1,\alpha,\ldots,\alpha^d$ ma $d+1$
elementów w przestrzeni wektorowej wymiaru $d$,
więc jest liniowo zależna. Istnieją
$a_0,\dots,a_d\in K$, nie wszystkie zerowe,
takie, że $\sum_{j=0}^d a_j\alpha^j=0$.
Oznacza to, że $\alpha$ zeruje niezerowy wielomian
stopnia co najwyżej $d$. Wielomian minimalny dzieli
ten wielomian, więc jego stopień nie przekracza $d$.
\end{proof}

\begin{wniosek}[Dołączenie skończenie wielu elementów algebraicznych]
\label{aa:wn:finite-algebraic-generators}
Niech $L/K$ będzie rozszerzeniem ciał, $r\ge1$,
a $\alpha_1,\ldots,\alpha_r\in L$ będą algebraiczne nad $K$.
Wówczas
$K(\alpha_1,\ldots,\alpha_r)/K$ jest rozszerzeniem
skończonym. Nie oznacza to, że każde rozszerzenie
algebraiczne ma skończony stopień.
\end{wniosek}
\begin{proof}
Niech $K_0=K$, $K_j=K(\alpha_1,\ldots,\alpha_j)$.
Wielomian o współczynnikach w $K$ zerujący
$\alpha_j$ jest także wielomianem nad $K_{j-1}$.
Zatem $\alpha_j$ jest algebraiczny nad $K_{j-1}$
i $[K_j:K_{j-1}]$ jest skończony.
Stosując kolejno twierdzenie o wieży dostajemy
$[K_r:K]=\prod_{j=1}^r[K_j:K_{j-1}]<\infty$.
Natomiast rozszerzenie
$E=\mathbb Q(2^{1/n}:n\ge1)\subseteq\mathbb R$
jest algebraiczne, lecz nieskończonego stopnia.
Istotnie, każdy jego element należy do ciała
$\mathbb Q(2^{1/n_1},\ldots,2^{1/n_k})$ dla
pewnej skończonej rodziny indeksów i jest algebraiczny
na mocy właśnie udowodnionego faktu. Z drugiej strony
dla każdego $n\ge1$ wielomian $X^n-2$ jest
nierozkładalny nad $\mathbb Q$ na mocy
Twierdzenia~\ref{aa:tw:eisenstein} (kryterium Eisensteina)
dla $p=2$: wszystkie współczynniki poza wiodącym
są podzielne przez $2$, a wyraz wolny $-2$
nie jest podzielny przez $4$. Stąd
$[\mathbb Q(2^{1/n}):\mathbb Q]=n$.
Gdyby $[E:\mathbb Q]$ było skończone, twierdzenie
o wieży ograniczałoby wszystkie te stopnie przez
$[E:\mathbb Q]$, co jest niemożliwe dla dowolnie
dużych $n$.
\end{proof}

\begin{twierdzenie}[Konstrukcja przez dołączenie formalnego pierwiastka]
\label{aa:tw:adjoin-root-abstract}
Niech $K$ będzie ciałem, a $f\in K[X]$ wielomianem
nierozkładalnym dodatniego stopnia. Wówczas
$E=K[X]/(f)$ jest ciałem zawierającym kanoniczny obraz
$K$, a $\alpha=[X]\in E$ spełnia $f(\alpha)=0$.
Ponadto $[E:K]=\deg f$. Dla każdego rozszerzenia
$L/K$ i pierwiastka $\beta\in L$ wielomianu $f$
istnieje dokładnie jedno $K$-zanurzenie
$E\to L$ wysyłające $\alpha$ na $\beta$.
Przez $K$-zanurzenie rozumiemy różnowartościowy
homomorfizm ciał, który na $K$ jest tożsamością.
\end{twierdzenie}
\begin{proof}
Nierozkładalność $f$ i algorytm Euklidesa w $K[X]$
pokazują, tak jak w dowodzie
Twierdzenia~\ref{aa:tw:simple-algebraic-extension},
że każdy niezerowy element ilorazu $K[X]/(f)$
ma odwrotność. Włączenie wielomianów stałych
$K\hookrightarrow K[X]/(f)$ jest różnowartościowe:
żaden niezerowy wielomian stały nie jest wielokrotnością
wielomianu dodatniego stopnia $f$.
Z definicji ilorazu $[f(X)]=0$, więc $f([X])=0$.
Jednoznaczność reszt stopnia mniejszego niż $\deg f$
daje bazę $1,\alpha,\dots,\alpha^{\deg f-1}$.

Jeżeli $f(\beta)=0$, to homomorfizm podstawiania
$\operatorname{ev}_\beta:K[X]\to L$ zeruje $(f)$.
Zatem przechodzi na homomorfizm
$\overline{\operatorname{ev}}_\beta:K[X]/(f)\to L$,
który wysyła $[X]$ na $\beta$ i pozostawia
współczynniki z $K$ bez zmian. Jest różnowartościowy,
ponieważ jego dziedzina jest ciałem, więc jądro jest
ideałem równym $0$ albo całemu ciału; drugi przypadek
wyklucza warunek $\overline{\operatorname{ev}}_\beta(1)=1$.
Każdy homomorfizm
o wskazanych wartościach musi wysyłać klasę
$\sum_j a_j[X]^j$ na $\sum_j a_j\beta^j$,
co daje jednoznaczność.
\end{proof}

\begin{przyklad}[Ciało czteroelementowe]
\label{aa:prz:four-element-field}
W $\mathbb F_2[X]$ wielomian $f=X^2+X+1$ nie ma
pierwiastków, bo $f(0)=f(1)=1$; dla stopnia dwa
oznacza to nierozkładalność. Definiujemy
$\mathbb F_4=\mathbb F_2[X]/(f)$ i
$\alpha=[X]$. Wtedy $\alpha^2=\alpha+1$,
ponieważ w charakterystyce dwa mamy $-1=1$.
Ciało zawiera dokładnie cztery elementy
$0,1,\alpha,1+\alpha$, a mnożenie wyznacza
relacja $\alpha^2=\alpha+1$.
W szczególności
$\alpha(\alpha+1)=\alpha^2+\alpha=1$,
więc $\alpha^{-1}=\alpha+1$. Co więcej,
$\alpha^3=\alpha(\alpha+1)=1$, zatem $\alpha$
ma rząd trzy w grupie multiplikatywnej, ponieważ
$\alpha\ne1$ i $3$ jest liczbą pierwszą.
\end{przyklad}

\section{Pierwiastki i ciała rozkładu}
\begin{definicja}[Ciało rozkładu]
\label{aa:def:splitting-field}
Niech $K$ będzie ciałem, a $f\in K[X]$ niezerowym wielomianem
dodatniego stopnia. Rozszerzenie $L/K$ nazywamy
\emph{ciałem rozkładu wielomianu $f$}, jeśli
$f$ rozkłada się w $L[X]$ na czynniki liniowe
(z uwzględnieniem jego niezerowego współczynnika wiodącego)
oraz $L$ jest generowane nad $K$ przez wszystkie
pierwiastki $f$ w $L$.
\end{definicja}

\begin{twierdzenie}[Istnienie ciała rozkładu]
\label{aa:tw:splitting-existence}
Dla każdego ciała $K$ i każdego niezerowego
$f\in K[X]$ dodatniego stopnia istnieje skończone
rozszerzenie $L/K$ będące ciałem rozkładu $f$.
\end{twierdzenie}
\begin{proof}
Przeprowadzimy indukcję po stopniu wielomianu.
Jeśli $\deg f=1$, ma on pierwiastek w $K$,
a jako ciało rozkładu możemy wziąć $L=K$.
Załóżmy, że twierdzenie jest prawdziwe dla
wszystkich wielomianów o stopniu mniejszym niż
$n$ i niech $\deg f=n\ge2$.
Wybierzmy nierozkładalny czynnik $g$ wielomianu $f$
w $K[X]$; taki czynnik istnieje, gdyż przy każdym
nietrywialnym rozkładzie stopnie czynników są mniejsze,
więc kolejne rozkładanie musi się zakończyć. Następnie
zbudujmy ciało
$E=K[X]/(g)$ z pierwiastkiem $\alpha=[X]$.
Twierdzenie~\ref{aa:tw:adjoin-root-abstract}
pokazuje, że $[E:K]=\deg g<\infty$.
W $E[X]$ możemy podzielić $f$ przez $X-\alpha$,
ponieważ $f(\alpha)=0$. Otrzymujemy
$f=(X-\alpha)h$ z $h\in E[X]$ i
$\deg h=n-1$.
Z założenia indukcyjnego istnieje skończone
rozszerzenie $L/E$, w którym $h$ rozkłada się
na czynniki liniowe i które jest generowane przez
pierwiastki $h$. Wtedy w $L$ rozkłada się również
$f$, a $L$ jest generowane nad $K$ przez
$\alpha$ oraz pierwiastki $h$, czyli przez
pierwiastki $f$. Na mocy twierdzenia o wieży
rozszerzenie $L/K$ jest skończone.
\end{proof}

\begin{przyklad}[Ciało rozkładu wielomianu sześciennego]
\label{aa:prz:cubic-splitting}
Wielomian $f=X^3-2$ jest nierozkładalny nad
$\mathbb Q$ na mocy Twierdzenia~\ref{aa:tw:eisenstein}
(kryterium Eisensteina) dla $p=2$.
Niech $\alpha=\sqrt[3]2\in\mathbb R$ i
$\omega=(-1+\sqrt{-3})/2\in\mathbb C$.
Bezpośredni rachunek daje $\omega^2+\omega+1=0$,
a $\omega\ne1$; z tożsamości
$(X-1)(X^2+X+1)=X^3-1$ wynika $\omega^3=1$.
Pierwiastkami $f$ są
$\alpha,\alpha\omega,\alpha\omega^2$.
Ciało rozkładu ma postać
$L=\mathbb Q(\alpha,\omega)$.
Mamy $[\mathbb Q(\alpha):\mathbb Q]=3$.
Wielomian $X^2+X+1$ nie ma pierwiastków rzeczywistych:
gdyby $x\in\mathbb R$ był jego pierwiastkiem, to
$(2x+1)^2=-3$, co jest niemożliwe. Ponieważ
$\mathbb Q(\alpha)\subseteq\mathbb R$, wielomian ten
nie ma pierwiastka również w $\mathbb Q(\alpha)$;
ma stopień dwa, więc jest nad nim nierozkładalny.
Stąd $[L:\mathbb Q(\alpha)]=2$, a twierdzenie
o wieży daje $[L:\mathbb Q]=6$.
Nie wystarcza samo dołączenie rzeczywistego
pierwiastka: $\mathbb Q(\alpha)$ nie zawiera
pozostałych dwóch pierwiastków.
\end{przyklad}

\section{Rozdzielczość wielomianów}
\begin{definicja}[Pierwiastek wielokrotny i rozdzielczość]
\label{aa:def:separable-polynomial}
Niech $K$ będzie ciałem i $f\in K[X]$ niezerowym
wielomianem dodatniego stopnia. Niech $r\ge1$ będzie
liczbą całkowitą. Pierwiastek $\alpha$
w pewnym rozszerzeniu $L/K$ ma \emph{krotność co najmniej
$r$}, jeśli $(X-\alpha)^r\mid f$ w $L[X]$.
Wielomian $f$ jest \emph{rozdzielczy}, jeśli ma
$\deg f$ różnych pierwiastków w swoim ciele rozkładu.
\end{definicja}

\begin{twierdzenie}[Kryterium rozdzielczości przez pochodną]
\label{aa:tw:separability-derivative}
Niech $K$ będzie ciałem, a $f=\sum_{j=0}^n a_jX^j\in K[X]$
niezerowym wielomianem stopnia $n\ge1$.
Wtedy $f$ jest rozdzielczy dokładnie wtedy, gdy
$\gcd(f,f')=1$ w $K[X]$, gdzie
$f'=\sum_{j=1}^n ja_jX^{j-1}$ jest formalną
pochodną wielomianu $f$.
\end{twierdzenie}
\begin{proof}
Dla $\alpha$ w rozszerzeniu $L/K$ zapiszmy
$f=(X-\alpha)q$. Reguła różniczkowania iloczynu
w pierścieniu wielomianów daje
$f'=q+(X-\alpha)q'$, więc $f'(\alpha)=q(\alpha)$.
Ponieważ $X-\alpha$ dzieli $q$ dokładnie wtedy,
gdy $q(\alpha)=0$, pierwiastek $\alpha$ ma krotność
co najmniej dwa wtedy i tylko wtedy, gdy
$f(\alpha)=f'(\alpha)=0$.

Jeśli $d=\gcd(f,f')$ ma dodatni stopień,
możemy wybrać jego nierozkładalny czynnik $h$.
W ciele rozkładu $h$ ma pierwiastek $\alpha$;
wtedy $f(\alpha)=f'(\alpha)=0$, więc $f$
ma pierwiastek wielokrotny i nie jest rozdzielczy.
Odwrotnie, jeśli $f$ ma pierwiastek wielokrotny
$\alpha$ w ciele rozkładu $L$, to oba wielomiany
$f,f'$ zerują $\alpha$. Wielomian minimalny
$m_{\alpha,K}$ dzieli więc jednocześnie $f$ i $f'$
na mocy Lematu~\ref{aa:lem:minimal-polynomial}.
Stąd $\deg\gcd(f,f')\ge1$. Dowiedliśmy równoważności.
\end{proof}

\begin{wniosek}[Nierozkładalność w charakterystyce zero]
\label{aa:wn:charzero-separable}
Każdy nierozkładalny wielomian dodatniego stopnia
nad ciałem charakterystyki zero jest rozdzielczy.
\end{wniosek}
\begin{proof}
Dla $f\in K[X]$ nierozkładalnego mamy
$f'\ne0$ w charakterystyce zero oraz
$\deg f'<\deg f$. Jedyny wspólny dzielnik
$f$ i $f'$ może być jednostką, bo każdy
niejednostkowy dzielnik $f$ jest stowarzyszony z $f$.
Zatem $\gcd(f,f')=1$ i poprzednie twierdzenie daje
rozdzielczość.
\end{proof}

\begin{przyklad}[Rozszerzenie nierozdzielcze]
\label{aa:prz:inseparable-example}
Niech $K=\mathbb F_p(t)$ i $f=X^p-t\in K[X]$.
Wielomian ten jest nierozkładalny nad $K$:
w $\mathbb F_p[t][X]$ spełnia założenia
Twierdzenia~\ref{aa:tw:eisenstein} (kryterium Eisensteina)
dla elementu pierwszego $t$, ponieważ $t$
dzieli współczynniki niewiodące, lecz $t^2\nmid -t$.
Wspomniane kryterium stwierdza nierozkładalność
właśnie nad ciałem ułamków $K$.
Jednak formalna pochodna $f'=pX^{p-1}=0$,
więc $\gcd(f,f')=f\ne1$. Wielomian jest
nierozdzielczy. Jeśli $\alpha^p=t$ w pewnym
rozszerzeniu, to w charakterystyce $p$ zachodzi
$X^p-t=(X-\alpha)^p$. Wielomian ma tylko jeden
różny pierwiastek w swoim ciele rozkładu,
i występuje on z krotnością $p$.
\end{przyklad}

\section{Ciała skończone}
\begin{twierdzenie}[Liczność ciała skończonego]
\label{aa:tw:finite-field-cardinality}
Każde ciało skończone $F$ ma $p^n$ elementów
dla pewnej liczby pierwszej $p$ i pewnej liczby
całkowitej $n\ge1$.
\end{twierdzenie}
\begin{proof}
Ciało $F$ ma dodatnią charakterystykę $p$,
ponieważ homomorfizm $\mathbb Z\to F$ nie może być
różnowartościowy: $\mathbb Z$ jest nieskończony,
a $F$ skończone. Z
Twierdzenia~\ref{aa:tw:prime-subfield} wynika,
że $F$ zawiera $\mathbb F_p$. Jako skończona
przestrzeń wektorowa nad tym podciałem ma pewien
wymiar $n\ge1$. Wybierając bazę, utożsamiamy
każdy element z jednoznaczną $n$-ką współczynników
z $\mathbb F_p$. Istnieje dokładnie $p^n$ takich
$n$-ek, zatem $|F|=p^n$.
\end{proof}

\begin{twierdzenie}[Grupa multiplikatywna ciała skończonego]
\label{aa:tw:finite-field-multiplicative-cyclic}
Dla każdego skończonego ciała $F$ grupa
$(F^\times,\cdot)$ jest cykliczna.
\end{twierdzenie}
\begin{proof}
Grupa $G=F^\times$ jest skończona i abelowa.
Jeśli $|G|=1$ (w szczególności dla $F=\mathbb F_2$),
to $G=\{1\}$ jest cykliczna. Dalej załóżmy $|G|>1$.
Na mocy Wniosku~\ref{aa:wn:abelian-fg-structure}
o klasyfikacji skończenie generowanych grup abelowych
jest sumą
prostą skończenie wielu grup cyklicznych.
Dla każdej liczby pierwszej $\ell$ dzielącej $|G|$
wybierzmy składnik cykliczny o największym rzędzie
$\ell^{a_\ell}$ w części $\ell$-pierwotnej
i generator $g_\ell$ tego składnika. W części
$\ell$-pierwotnej rząd każdego elementu dzieli
$\ell^{a_\ell}$, gdyż każdy jej składnik cykliczny
ma rząd dzielący $\ell^{a_\ell}$. Dlatego
wykładnik grupy $G$ (najmniejsza dodatnia liczba $e$
spełniająca $x^e=1$ dla wszystkich $x\in G$) jest
równy $e=\prod_\ell\ell^{a_\ell}$: liczba ta
anihiluje wszystkie składniki, a musi dzielić się
przez rząd każdego $g_\ell$. Element
$g=\prod_\ell g_\ell$ ma rząd równy $e$, gdyż
rzędy jego czynników są parami względnie pierwsze:
z $g^k=1$ wynika, po rzutowaniu na każdą część
pierwotną, $g_\ell^k=1$, więc
$\ell^{a_\ell}\mid k$ dla wszystkich $\ell$.
Odwrotnie, $g^e=1$.

Każdy element $G$ jest pierwiastkiem wielomianu
$X^e-1\in F[X]$. Niezerowy wielomian stopnia $e$
nad ciałem ma co najwyżej $e$ różnych pierwiastków,
zatem $|G|\le e$. Z drugiej strony $e$ jest
rzędem pewnego elementu $G$, więc z twierdzenia
Lagrange'a wynika $e\le|G|$. Stąd $e=|G|$,
a wybrany element rzędu $e$ generuje całe $G$.
\end{proof}

\begin{przyklad}[Dlaczego $\mathbb Z/4\mathbb Z$ nie jest ciałem]
\label{aa:prz:not-field-four}
Pierścień $\mathbb Z/4\mathbb Z$ ma cztery elementy,
ale $[2]^2=[0]$ i $[2]\ne[0]$, więc ma dzielniki
zera. Nie jest ciałem, mimo że istnieje
czteroelementowe ciało $\mathbb F_4$ skonstruowane
w Przykładzie~\ref{aa:prz:four-element-field}.
Równa liczba elementów nie oznacza izomorfizmu
struktur algebraicznych.
\end{przyklad}

\begin{uwaga}[Kolejna część: teoria Galois]
\label{aa:uw:field-extensions-future}
Skonstruowaliśmy rozszerzenia przez pierwiastki,
ciała rozkładu i poznaliśmy rolę rozdzielczości.
W następnych rozdziałach zbadamy przedłużanie
zanurzeń, automorfizmy rozszerzeń, normalność
oraz podstawowe twierdzenie teorii Galois.
Do wyjaśnienia, które równania można rozwiązać
przez pierwiastniki, potrzebne będą ponadto grupy
rozwiązalne z wcześniejszej części skryptu.
\end{uwaga}
